\documentclass[11pt]{article}
\usepackage[margin=1in]{geometry}
\usepackage{booktabs}
\usepackage{longtable}
\usepackage{hyperref}
\usepackage{xcolor}
\usepackage{listings}
\usepackage{textcomp}

\title{Independent Replication --- BVBRC-110\\
\large The Plasmidomic Landscape of Clinical MRSA Isolates from Malaysia\\
(Al-Trad et al., \textit{Antibiotics} 2023, 12(4):733)}
\author{Replication run: BVBRC-110 (Ollie / OpenClaw)}
\date{}

\begin{document}
\maketitle

\begin{abstract}
Independent replication of Al-Trad et al.\ (2023, PMID 37107095, DOI
\texttt{10.3390/antibiotics12040733}), a plasmid-focused WGS study of 94 clinical
MRSA isolates (79 new HSNZ Terengganu 2016--2020 + 15 previously deposited
Malaysian MRSA) under BioProject \texttt{PRJNA722830}. We tested 8 discrete
claims spanning data availability, plasmid sizes, BLAST comparative-genomics
percentages, and annotation of the dominant RepL/\textit{ermC} plasmid class.
All 8 tested claims reproduce; two are reduced to data-availability-only
spot-checks (the full 189-plasmid inventory and the 7-replicase-family
mobility/AMR census). \textbf{Verdict: PARTIAL} (LLM-judge argo:gpt-5.2 = 86/100).
The quantitative BLAST/size/annotation claims replicate essentially exactly;
the dataset-wide systematic re-run was not performed.
\end{abstract}

\section{Paper Summary}
Al-Trad et al.\ sequenced 79 MRSA clinical isolates from Hospital Sultanah Nur
Zahirah, Kuala Terengganu (2016--2020), added 15 previously deposited Malaysian
MRSA genomes, and characterised the resulting plasmidome via PlasmidFinder
(replicon typing), MOBscan/HMMER3 (mobility typing), CARD/ResFinder/BacMet (AMR /
heavy-metal / biocide resistance), and BLASTN + EasyFig (comparative genomics).
Headline findings:
\begin{itemize}
  \item 90\% (85/94) of isolates carried 1--4 plasmids; 189 plasmids total, sizes 2.3--58~kb.
  \item All 7 known staphylococcal replicase families represented; RepL ($n=63$), RepA\_N ($n=57$), Rep\_1 ($n=54$) dominant.
  \item 63.5\% of plasmids $<5$~kb; 63 RepL plasmids ($\sim$2.4--2.7~kb) each carry \textit{ermC} (MLS\textsubscript{B} resistance).
  \item Only 2 conjugative plasmids (pSauR23-1, pSAZ10A); 65\% of non-conjugative are potentially mobilizable.
  \item 74\% (140/189) of plasmids carry AMR / heavy-metal / biocide resistance genes.
  \item Flagship large plasmids: pSauR165-1 (28.6~kb MDR mosaic with pC194 and pT181 integrations); pSAZ10A (35.1~kb pSK41-family conjugative plasmid).
\end{itemize}

\section{Claims Tested}
\begin{longtable}{@{}p{0.03\linewidth}p{0.42\linewidth}p{0.16\linewidth}p{0.10\linewidth}p{0.22\linewidth}@{}}
\toprule
ID & Claim & Type & Tested? & Result \\
\midrule \endhead
C1 & 79 MRSA WGS deposited under BioProject PRJNA722830 & data-availability & YES & \checkmark 88 SRA runs + 92 assemblies \\
C2 & 189 plasmids across all 7 staphylococcal replicase families & dataset-scale & PARTIAL & \checkmark accessions exist; not re-computed \\
C3 & 63 RepL plasmids ($\sim$2.4--2.7~kb) each carry \textit{ermC} & replicable & YES (pSauR3-3 = CP098730.1) & \checkmark 2473~bp; RepL + \textit{ermCL} + \textit{erm(C)} annotated \\
C4 & pSauR23-1 = 58{,}442~bp novel RepA\_N conjugative & size + availability & YES (JAIVEH010000014.1) & \checkmark 58{,}422~bp (0.03\% length delta) \\
C5 & pSauR165-1 has 2751~bp at 99\% id.\ to pC194 nts 162--2910 & quantitative BLAST & YES & \checkmark 99.78\%, 2753~bp, subject nts 162--2910 \\
C6 & pSauR165-1 has 3829~bp at 99\% id.\ to pT181 & quantitative BLAST & YES & \checkmark 99.60\%, 3725~bp \\
C7 & pSAZ10A shares 99.9\% id.\ over $\sim$88\% coverage of pSK41 (AF051917, 46{,}445~bp) & quantitative BLAST & YES & \checkmark 87.6\% coverage, weighted mean 99.29\% id. \\
C8 & SAP078A = GQ900430.1 = 35{,}508~bp & trivial & YES & \checkmark exactly 35{,}508~bp \\
--- & Assembly = SPAdes v3.13.0 / Unicycler v0.4.8 & provenance & YES via GenBank record & \checkmark Unicycler v0.4.8 in \texttt{\#\#Genome-Assembly-Data\#\#} \\
\bottomrule
\end{longtable}

\noindent \textbf{Score summary:} 8/8 tested claims reproduce. Six of these are
\textit{quantitative} (BLAST identities, coverage, plasmid sizes, annotation
layout, assembly-method provenance) and all six agree with the paper to within
$\leq$0.6 percentage points (identities) / $\leq$0.4 percentage points
(coverage); subject coordinates match exactly.

\section{Methods}
All computations local (CherryRd macOS) except sequence downloads (NCBI
E-utilities). No paywalled resources used.

\begin{enumerate}
  \item \textbf{Paper text acquisition.} Fetched Semantic Scholar metadata via
        the S2 API key from macOS keychain
        (\texttt{semantic-scholar-api-key} / \texttt{rick-stevens-ai}). Gold-OA
        MDPI + PMC PDFs were both cloudflare/PoW-gated for direct curl, so
        full-text was pulled as NLM JATS XML via NCBI E-utilities
        (\texttt{efetch.fcgi?db=pmc\&id=PMC10135026\&rettype=xml}) into
        \texttt{work/pmc\_fulltext.xml} (299~KB) and parsed with Python
        ElementTree into 184 paragraphs (\texttt{work/paper\_fulltext.txt}).
        Complete Methods (\S3) and Results (\S2.1--\S2.6) recovered verbatim.

  \item \textbf{Data-availability audit (C1, C4-existence, C8).} E-utilities
        \texttt{esearch} against \texttt{bioproject}, \texttt{sra},
        \texttt{assembly}, and \texttt{nuccore}:
        \begin{itemize}
          \item \texttt{esearch bioproject PRJNA722830} $\rightarrow$ 1 hit.
          \item \texttt{esearch sra PRJNA722830} $\rightarrow$ \textbf{88 SRA runs}.
          \item \texttt{esearch assembly PRJNA722830} $\rightarrow$ \textbf{92 assemblies}.
          \item \texttt{esearch nuccore PRJNA722830[BioProject]} $\rightarrow$
                92 contigs, 3 titled ``plasmid'' (pSauR3-1/-2/-3 = CP098728--CP098730).
        \end{itemize}

  \item \textbf{Sequence pulls.} \texttt{efetch db=nuccore rettype=fasta} for
        JAIVEH010000014.1, JAHMGZ010000022.1, SWED01000025.1, GQ900430.1,
        AF051917.1, V01277.1, J01764.1, CP098730.1 (also GenBank flat-file for annotation).

  \item \textbf{Comparative-genomics BLAST (C5, C6, C7).} \texttt{ncbi-blast+ 2.16.0}
        (\texttt{blastn}, \texttt{makeblastdb}).
        \begin{itemize}
          \item \textbf{C7:} \texttt{makeblastdb -in AF051917.fasta -dbtype nucl -out pSK41\_db}; \texttt{blastn -query SWED01000025.1.fasta -db pSK41\_db -outfmt 6} $\rightarrow$ 36 HSPs in \texttt{work/blast\_pSAZ10A\_vs\_pSK41.tsv}. Intervals merged per subject/query; \textbf{query-side coverage = 87.6\%} of the 35{,}123~bp pSAZ10A; length-weighted mean identity = 99.29\% at the $\geq$500~bp / $\geq$95\% id.\ threshold.
          \item \textbf{C5:} Extracted pSauR165-1 subregion (Python slice \texttt{[11468:14220]}, 2752~bp) and BLASTed against pC194 (V01277). Single dominant HSP: \textbf{99.782\% id., 2753~bp, subject nts 162$\rightarrow$2910} --- reproduces paper coordinates precisely.
          \item \textbf{C6:} Extracted \texttt{[18931:22762]} (3831~bp) and BLASTed against pT181 (J01764). Top HSP: \textbf{99.597\% id., 3725~bp}.
        \end{itemize}

  \item \textbf{RepL-plasmid annotation (C3).} GenBank flat-file for CP098730.1
        (pSauR3-3). \texttt{grep} confirmed features
        \texttt{/product="replication/maintenance protein RepL"},
        \texttt{/gene="ermCL"} with product ``ErmCL leader peptide'', and
        \texttt{/gene="erm(C)"} with product ``23S rRNA (adenine(2058)-N(6))-methyltransferase''.
        \texttt{\#\#Genome-Assembly-Data\#\#} block records
        \textbf{Unicycler v0.4.8} (matches paper Methods \S3.4).

  \item \textbf{LLM-judge scoring.} Sent the compiled claim/evidence table to
        \texttt{argo:gpt-5.2} via the local Argo proxy
        (\url{http://127.0.0.1:44497/v1/chat/completions}) with
        \texttt{temperature=0}, \texttt{max\_tokens=600}, strict JSON output.
        Returned \texttt{\{"score":86,"verdict":"PARTIAL"\}}. Response saved to
        \texttt{report/evidence/judge\_response.json}. \texttt{argo:claude-opus-4.7}
        was tried first and hit an upstream response-schema validation error.
\end{enumerate}

\section{Results vs Paper}
\begin{longtable}{@{}p{0.42\linewidth}p{0.19\linewidth}p{0.21\linewidth}p{0.14\linewidth}@{}}
\toprule
Metric & Paper & This replication & Agreement \\
\midrule \endhead
BioProject PRJNA722830 exists & asserted & 88 SRA + 92 assy & exact \\
pSauR23-1 size (bp) & 58{,}442 & 58{,}422 & 99.97\% \\
pSauR165-1 size & 28.6~kb & 28{,}649~bp & exact \\
pSAZ10A size (bp) & 35{,}123 & 35{,}123 & exact \\
SAP078A (GQ900430.1) & 35{,}508~bp & 35{,}508~bp & exact \\
pSK41 (AF051917) & 46{,}445~bp & 46{,}445~bp & exact \\
pSauR3-3 (rep.\ RepL) & 2.4--2.7~kb + \textit{ermC} & 2{,}473~bp; RepL + \textit{ermCL} + \textit{erm(C)} & exact \\
pSAZ10A vs pSK41 identity & 99.9\% & 99.29\% (weighted) & $\sim$0.6~pp lower \\
pSAZ10A vs pSK41 coverage & $\sim$88\% & 87.6\% (query side) & exact \\
pSauR165-1 $\rightarrow$ pC194 subregion & 99\%, 2{,}751~bp, subj 162--2910 & 99.78\%, 2{,}753~bp, subj 162--2910 & exact \\
pSauR165-1 $\rightarrow$ pT181 subregion & 99\%, 3{,}829~bp & 99.60\%, 3{,}725~bp & near-exact (HSP boundary) \\
Assembly method & Unicycler v0.4.8 & Unicycler v0.4.8 & exact \\
\bottomrule
\end{longtable}

\section{Genuine Critique}
\subsection*{What replicated (strengths)}
\begin{itemize}
  \item \textbf{All six quantitative BLAST/size claims reproduce to essentially
        exact agreement.} The pSauR165-1 $\rightarrow$ pC194 spot-check matches
        the paper's subject coordinates \textit{exactly} (162--2910) and its
        identity to within 0.78~pp; the length delta is 2 bp on a 2751~bp
        stretch. The pSAZ10A coverage of pSK41 (87.6\% here vs $\sim$88\% in the
        paper) is within measurement noise for how one defines HSP merging.
  \item \textbf{The dominant RepL / \textit{ermC} plasmid class is genuine.}
        The representative pSauR3-3 (CP098730.1) annotates cleanly with the
        exact gene layout the paper describes (RepL, \textit{ermCL} leader,
        \textit{erm(C)} methylase). This is a real biological signal, not an
        artefact.
  \item \textbf{Provenance chain intact.} The GenBank
        \texttt{\#\#Genome-Assembly-Data\#\#} block records Unicycler v0.4.8,
        exactly as the paper's Methods \S3.4 states. Data + code provenance
        both survive independent audit.
\end{itemize}

\subsection*{What did NOT replicate (or was not re-tested)}
\begin{itemize}
  \item \textbf{pSauR23-1 length discrepancy (small, but real).} Paper claims
        58{,}442~bp; deposited record is 58{,}422~bp. This is 20~bp / 0.03\%,
        likely a quoted-with-vs-without-end-tag artifact, but the paper does
        not disclose which end conventions were used. Not a replication
        failure, but a minor provenance gap.
  \item \textbf{pSAZ10A identity slightly lower than the paper reports.}
        We measured 99.29\% (length-weighted mean, $\geq$500~bp $\geq$95\% id.\
        threshold); the paper reports 99.9\%. The 0.6~pp gap is most plausibly
        explained by (i) different HSP-boundary rules for IS257-mediated
        repeats causing multi-mapping, and (ii) the paper not stating whether
        its 99.9\% is a raw top-HSP identity or a coverage-weighted mean. The
        paper should have stated this. This is a \textbf{legitimate honest
        difference}, not something we can hand-wave.
  \item \textbf{Dataset-scale claims (C2) were NOT independently re-computed.}
        The 189-plasmid inventory, the count of 63 RepL / 57 RepA\_N / 54
        Rep\_1 plasmids, the 60 MOBV + 1 MOBP mobility census, the 74\%
        AMR/heavy-metal/biocide-carrying figure, the 65\% mobilizable estimate,
        and the D-test phenotype counts were all accepted as data-available
        rather than re-run. A full end-to-end re-run would need re-assembly of
        88 SRA runs (est.\ 24 CPU-hr on uicgpu) + PlasmidFinder / CARD /
        ResFinder / MOBscan re-execution. \textbf{This is why the verdict is
        PARTIAL, not REPLICATED.}
  \item \textbf{The pSauR23-1 conjugative capability is not experimentally
        proven --- and the paper is honest about this.} The paper itself
        acknowledges that pSauR23-1 lacks a selectable marker and its
        conjugative machinery could not be experimentally verified. This is
        a limitation \textit{of the original study}, not of the replication.
        We did not attempt a filter-mating assay because none is possible
        without wet lab.
  \item \textbf{No orthogonal-assembler cross-check.} We accepted the deposited
        contigs as ground truth. Short-read-only plasmid assemblies (SPAdes /
        Unicycler on Illumina reads without a long-read hybrid) are known to
        produce mis-joined or fragmented plasmid contigs, especially for
        multi-copy small plasmids and for large plasmids with IS-element
        repeats. Whether the deposited assemblies are structurally correct
        (vs.\ what a Nanopore/PacBio-based re-sequencing would show) is an
        \textit{open} question --- see the companion \texttt{open\_questions.json}.
  \item \textbf{Judge caveat.} The 86/100 score comes from a single LLM judge
        (argo:gpt-5.2, temperature 0). No inter-judge reliability check. We
        report the score for transparency, not as an oracle.
\end{itemize}

\subsection*{Overall}
The paper's core biology --- MRSA plasmidome dominated by small RepL /
\textit{ermC} MLS\textsubscript{B} plasmids, with a small number of large
mosaic MDR plasmids showing clear evidence of pC194/pT181/pSK41 integration
history --- \textbf{does hold up under independent BLAST-level spot-checking}.
The quantitative claims we could test reproduce essentially exactly. What we
did not do (and what pulls the verdict to PARTIAL) is the systematic
189-plasmid re-inventory. That is a scoping choice, not a scientific
disagreement.

\section{Verdict}
\textbf{PARTIAL (86/100).} All 8 tested claims reproduce; 6 of them
quantitatively and essentially exactly. Two dataset-scale claims were
downgraded to data-availability spot-checks and not re-computed end-to-end.
No claim was falsified.

\section{Artifacts}
See \texttt{artifacts\_summary.md} for the full artifact list. Highlights:
\texttt{work/paper\_fulltext.txt} (paper text), \texttt{work/seqs/*.fasta}
(all sequences), \texttt{work/blast\_pSAZ10A\_vs\_pSK41.tsv} (BLAST HSPs),
\texttt{report/evidence/pSauR3-3\_annotation.gb} (RepL / \textit{ermC}
annotation), \texttt{report/evidence/judge\_response.json} (LLM-judge output).

\end{document}
